% Supplementary Table: robustness of the 4N trajectory across the ten
% lowest-objective separate in vitro fits.
%
% PURPOSE
% This table documents the per-fit values underlying the Results statement
% that the modeled 4N chromosome-number decline was 3.54--4.17 chromosomes
% and that direct hypoxia-origin dead burden remained below 1.9% across the
% ten lowest-objective separate in vitro fits.
%
% PRIMARY RESULT ROOT
%   /Users/4482173/Documents/GitHub/soft_coupling/oxygen/results/
%   fit_invitro_O2_buffering_500seed/
%
% PER-SEED SOURCE FILES
% 1. Objective:
%    seed*/fit_summary.tsv
%    row: metric = objective_total
% 2. Initial and terminal modeled mean chromosome number:
%    seed*/invitro_lineage_summary.tsv
% 3. Direct hypoxia-origin dead burden and total burden:
%    seed*/invitro_daily_counts.tsv
%
% TOP-TEN SELECTION AND PORTABLE AUDIT FILES
% 1. Ranked source-directory index:
%    /Users/4482173/Documents/GitHub/HypoxiaLTEEFigures/revised/SuppFiles/
%    round1/miningcloneid_review_v2/analysis_tables/
%    source_top10_index_localized.csv
%    Selection: result_group = fit_invitro_O2_buffering_500seed and
%    rank = 1,...,10 after ascending objective ranking.
% 2. Exact per-seed values used in this table:
%    /Users/4482173/Documents/GitHub/HypoxiaLTEEFigures/revised/SuppFiles/
%    round1/miningcloneid_review_v2/analysis/
%    invitro_top10_robustness.csv
% 3. Recalculation script:
%    /Users/4482173/Documents/GitHub/HypoxiaLTEEFigures/revised/SuppFiles/
%    round1/miningcloneid_review_v2/scripts/reanalyse_results.py
%    Relevant block: "# ---------------- in vitro ----------------".
%
% CHROMOSOME-NUMBER CALCULATION
% For each selected seed:
% 1. Read invitro_lineage_summary.tsv.
% 2. Retain cohort = 4N.
% 3. Collapse repeated observation records sharing the same segment_id,
%    because parallel O1/O2 records at a matched segment share one modeled
%    lineage state.
% 4. Initial mean N = predicted_mean_kary_N at the minimum passage_index.
% 5. Terminal mean N = predicted_mean_kary_N at the maximum passage_index.
% 6. Decline in mean N = initial mean N - terminal mean N.
% The final modeled passage is on the terminal 0% O2 deprivation branch.
%
% DIRECT HYPOXIA-ORIGIN DEAD-BURDEN CALCULATION
% For each selected seed:
% 1. Read invitro_daily_counts.tsv and retain cohort = 4N.
% 2. At every row with burden_total > 0, calculate:
%      direct hypoxia-origin dead fraction
%        = burden_dead_hypoxia / burden_total.
% 3. Report the maximum of this fraction over all daily 4N states.
% In all ten selected fits, the maximum occurred at passage_index = 29,
% oxygen_pct = 0, day = 10, on the severe-deprivation branch.
% The displayed percentage is 100 times the fraction.
%
% IMPORTANT DENOMINATOR/TIMEPOINT NOTE
% This maximum-over-daily-states definition is the source of the 1.74--1.82%
% range. It is not the same summary as the passage-end fraction shown for
% seed10 in Figure 3F, which uses a selected passage endpoint. Both summaries
% use burden_total as the denominator, but they differ in timepoint selection.
%
% VERIFIED UNROUNDED RANGES
% Objective: 3.85253526260594--3.86230239787459.
% Initial mean N: 84.2083996718758--84.7084758541827.
% Terminal mean N: 80.1277453011119--80.8905850002582.
% Decline: 3.54463845865141--4.17236561451030 chromosomes.
% Maximum direct hypoxia-origin dead fraction:
%   0.0173926520856655--0.0181541247233420
%   = 1.73926520856655%--1.81541247233420%.
%
% INTERPRETIVE LIMIT
% The ten rows are optimizer-selected numerical fits, not ten biological
% replicates and not a confidence interval. The table establishes numerical
% robustness across retained near-optimal solutions; it is not an additional
% hypothesis test.

\begin{table}[!htbp]
\centering
\small
\setlength{\tabcolsep}{4.5pt}
\caption{Modeled 4N chromosome-number decline and direct hypoxia-origin dead burden across the ten lowest-objective separate \textit{in vitro} fits. Initial and terminal mean chromosome numbers refer to the first and last modeled states of the selected 4N lineage; its terminal portion was maintained at target 0\% O$_2$. The last column is the maximum daily direct hypoxia-origin dead burden divided by total burden over the full modeled 4N trajectory.}
\label{tab:invitro_top10_4n_trajectory}
\begin{tabular}{@{}rlrrrrr@{}}
\toprule
Rank & Fit & Objective & Initial mean \(N\) & Terminal mean \(N\) & Decline in mean \(N\) & \shortstack{Maximum direct hypoxia-origin\\dead burden (\%)} \\
\midrule
1  & \texttt{seed10}  & 3.8525 & 84.34 & 80.76 & 3.58 & 1.76 \\
2  & \texttt{seed132} & 3.8533 & 84.21 & 80.66 & 3.55 & 1.76 \\
3  & \texttt{seed81}  & 3.8541 & 84.22 & 80.67 & 3.54 & 1.77 \\
4  & \texttt{seed294} & 3.8594 & 84.71 & 80.89 & 3.82 & 1.74 \\
5  & \texttt{seed337} & 3.8598 & 84.66 & 80.76 & 3.91 & 1.74 \\
6  & \texttt{seed106} & 3.8605 & 84.45 & 80.53 & 3.91 & 1.82 \\
7  & \texttt{seed317} & 3.8610 & 84.30 & 80.13 & 4.17 & 1.76 \\
8  & \texttt{seed140} & 3.8610 & 84.21 & 80.54 & 3.67 & 1.74 \\
9  & \texttt{seed285} & 3.8618 & 84.24 & 80.20 & 4.04 & 1.77 \\
10 & \texttt{seed464} & 3.8623 & 84.68 & 80.73 & 3.96 & 1.77 \\
\bottomrule
\end{tabular}
\par\vspace{0.5em}
\begin{minipage}{0.96\textwidth}
\footnotesize
\textit{Note:} Fits are ordered by the retained objective value. The decline is the initial minus terminal predicted mean chromosome number. The dead-burden percentage is the maximum, over all daily 4N states, of \(100\times\texttt{burden\_dead\_hypoxia}/\texttt{burden\_total}\); it is therefore a trajectory maximum rather than an endpoint-only value. These optimizer-selected solutions are not biological replicates or a confidence interval.
\end{minipage}
\end{table}
